\chapter{Comparing Mutual Information with Normal Wald}
\label{chap:baselines}

This chapter defines the estimand and statistic shared by the two methods in
the final experiment. Normal Wald evaluates it against a standard normal
reference; Chapter~\ref{chap:expanded} develops a Student reference with
estimated degrees of freedom.

\section{Data, estimand, and assumptions}

Let
\begin{equation}
  N^P=(N^P_{ij}),
  \qquad
  N^Q=(N^Q_{ij})
\end{equation}
be independent \(r\times c\) count tables with totals \(n_P\) and \(n_Q\).
The row and column categories have the same definitions in both populations.
The empirical probabilities are
\begin{equation}
  \widehat p_{ij}=\frac{N^P_{ij}}{n_P},
  \qquad
  \widehat q_{ij}=\frac{N^Q_{ij}}{n_Q}.
\end{equation}

The population parameter is the signed difference
\begin{equation}
  \Delta=I(P)-I(Q).
  \label{eq:delta-population}
\end{equation}
The principal test is two-sided:
\begin{equation}
  H_0:\Delta=0
  \quad\text{against}\quad
  H_1:\Delta\ne0.
  \label{eq:two-sample-null}
\end{equation}

The first-order construction uses the following assumptions.

\begin{assumption}[Sampling]
Observations are independent and identically distributed within each
population, and the two population samples are independent.
\end{assumption}

\begin{assumption}[Aligned fixed alphabets]
The row and column categories are fixed in advance and aligned between the two
tables.
\end{assumption}

\begin{assumption}[Positive support]
Every population cell has positive probability. Empirical zero cells are
allowed, but widespread observed support loss is recorded as a diagnostic.
\end{assumption}

\begin{assumption}[Regular MI]
The first-order influence variances \(V(P)\) and \(V(Q)\) are positive. This
excludes exact independence and other degenerate cases from the regular
Student approximation.
\end{assumption}

All logarithms are natural. MI differences and standard errors are therefore
reported in nats. Changing the logarithm base scales both by the same positive
constant and leaves the standardised statistic unchanged.

\section{Plug-in and bias-corrected MI}

For the first table, define empirical pointwise MI by
\begin{equation}
  \widehat\pmi_P(i,j)
  =\log\left(
  \frac{\widehat p_{ij}}
       {\widehat p_{i+}\widehat p_{+j}}
  \right)
  \qquad\text{when }\widehat p_{ij}>0.
  \label{eq:empirical-pmi}
\end{equation}
The plug-in MI estimate is
\begin{equation}
  \Ihat(P)
  =\sum_{i,j:\widehat p_{ij}>0}
  \widehat p_{ij}\widehat\pmi_P(i,j).
  \label{eq:plugin-mi}
\end{equation}
The same calculation gives \(\Ihat(Q)\).

Let
\begin{equation}
  d=(r-1)(c-1).
\end{equation}
Under fixed positive support, the leading bias correction gives
\begin{equation}
  \widehat I_{\mathrm{BC}}(P)
  =\Ihat(P)-\frac{d}{2n_P},
  \qquad
  \widehat I_{\mathrm{BC}}(Q)
  =\Ihat(Q)-\frac{d}{2n_Q}.
  \label{eq:bias-corrected-mi}
\end{equation}
The estimated effect is
\begin{equation}
  \Deltahat
  =\widehat I_{\mathrm{BC}}(P)
  -\widehat I_{\mathrm{BC}}(Q).
  \label{eq:bias-corrected-difference}
\end{equation}

The correction uses the declared alphabet dimensions rather than the observed
nonempty dimensions. It may produce a negative corrected MI estimate in a
sparse sample. The implementation does not clip that estimate at zero. Clipping
would replace the stated bias-corrected estimator by a different nonlinear
estimator and would change its finite-sample bias and sampling distribution.
When $n_P=n_Q$, the two leading corrections cancel in the difference in
Equation~\eqref{eq:bias-corrected-difference}. When the sample sizes differ,
their difference remains and removes the leading unequal-sample bias described
in Section~\ref{sec:background-mi-bias}.

\section{First-order variance of plug-in MI}

For population \(P\), define
\begin{equation}
  V(P)
  =\sum_{i,j}p_{ij}
  \{\pmi_P(i,j)-I(P)\}^2.
  \label{eq:population-mi-variance}
\end{equation}
This is the population variance of pointwise MI. It governs the sampling
variance of the complete plug-in estimator because the first-order MI error is
an average of centred pointwise-MI contributions. This first-order variance
is established for histogram and discrete MI \citep{moddemeijer1989,brillinger2004};
the expanded calculation in Chapter~\ref{chap:expanded} concerns the
\emph{sampling uncertainty of its plug-in estimate}.

To see this directly, write one observation as
\(Z_a^{(P)}=(X_a^{(P)},Y_a^{(P)})\). A first-order Taylor expansion of MI
around the fixed population table $P$ gives
\begin{equation}
  \Ihat(P)-I(P)
  =
  \frac{1}{n_P}
  \sum_{a=1}^{n_P}
  \{\pmi_P(Z_a^{(P)})-I(P)\}
  +o_p(n_P^{-1/2}).
  \label{eq:mi-asymptotic-linear}
\end{equation}
The summands are independent, have mean zero, and have variance \(V(P)\).
The remainder $o_p(n_P^{-1/2})$ is asymptotically smaller than the leading
$n_P^{-1/2}$ sampling fluctuation. Appendix~\ref{app:derivation-details}
derives the linear term by expanding MI as a function of the cell
probabilities.
Therefore,
\begin{align}
  \Var\{\Ihat(P)\}
  &\approx
  \Var\left[
  \frac{1}{n_P}
  \sum_{a=1}^{n_P}
  \{\pmi_P(Z_a^{(P)})-I(P)\}
  \right]\\
  &=\frac{1}{n_P^2}
  \sum_{a=1}^{n_P}V(P)\\
  &=\frac{V(P)}{n_P}.
  \label{eq:mi-sampling-variance}
\end{align}
Subtracting the fixed leading bias term in Equation
\eqref{eq:bias-corrected-mi} does not change this first-order variance.

The observed table estimates \(V(P)\) by
\begin{equation}
  \Vhat(P)
  =\sum_{i,j:\widehat p_{ij}>0}
  \widehat p_{ij}
  \{\widehat\pmi_P(i,j)-\Ihat(P)\}^2.
  \label{eq:estimated-mi-variance}
\end{equation}
Define \(\Vhat(Q)\) analogously.

\section{Shared standard error and test statistic}

Because the two samples are independent, their estimated variance
contributions add:
\begin{equation}
  \SEhat^2
  =\frac{\Vhat(P)}{n_P}
  +\frac{\Vhat(Q)}{n_Q}.
  \label{eq:shared-se}
\end{equation}
The common standardised statistic is
\begin{equation}
  T=\frac{\Deltahat}{\SEhat}.
  \label{eq:shared-t}
\end{equation}
The null hypothesis is not imposed when calculating \(\Deltahat\),
\(\SEhat\) or \(T\): they use the observed tables. Equal population MI
enters when interpreting the reference distribution, which is centred at
zero under \(H_0\). In particular, \(Q\)'s observed estimate and variance
are not set to zero merely because the null is being tested.
Equations~\eqref{eq:bias-corrected-difference}--\eqref{eq:shared-t} are shared
by normal Wald and expanded Welch. Their difference lies in the reference
distribution, not the effect estimate or plug-in standard error.

\section{Normal Wald reference}

Under the regular first-order asymptotics,
\begin{equation}
  T\xrightarrow{d}N(0,1)
  \qquad\text{under }H_0.
\end{equation}
Here both $n_P$ and $n_Q$ increase while the aligned alphabet and the two
positive-support population tables remain fixed.
The two-sided normal Wald p-value is
\begin{equation}
  p_{\mathrm{Wald}}
  =2\{1-\Phi(|T|)\},
  \label{eq:wald-pvalue}
\end{equation}
where \(\Phi\) is the standard normal distribution function. A
\(100(1-\alpha)\%\) confidence interval for \(\Delta\) is
\begin{equation}
  \Deltahat
  \ \pm\
  z_{1-\alpha/2}\SEhat,
  \label{eq:wald-ci}
\end{equation}
where $z_\gamma$ is the $\gamma$ quantile of the standard normal
distribution.

The normal reference is a large-sample approximation. Its standard error is
estimated from the observed tables; Wald does not treat it as known. Rather,
it does not apply the finite-sample degrees-of-freedom adjustment proposed
here.

\section{What the expanded method must calculate}

The shared statistic supplies the estimated effect and its first-order
standard error. A familiar Welch analogy would assign \(n_P-1\) and \(n_Q-1\)
degrees of freedom, as for ordinary sample variances. But \(\Vhat(P)\) is a
nonlinear functional of a table whose margins and pointwise MI values also
change with the data. The expanded method instead estimates component degrees
of freedom from that full functional.

Satterthwaite moment matching requires the first two moments of each positive
variance estimate. To first order,
\begin{equation}
  \E\{\Vhat(P)\}\approx V(P).
\end{equation}
Chapter~\ref{chap:expanded} derives
\(\Var\{\Vhat(P)\}\) by calculating how the complete variance functional
changes when one observation changes the joint table.
